\subsection{Эквивалентность $\bigvee \limits_{i} D_{n_i, k_i}(x,y)$ и $D_{n,k}(x,y)$}

\begin{namedLemma}[Лемма о классах достижимости конечных объединений]
    \label{lemm:RCtoR}
    Пусть
    $$R_j(x,y)=\bigvee\limits_{i=1}^{N_j}D_{n_{j,i},k_{j,i}}(x,y),
    \qquad j\in\{1,2\},$$
    и каждое из отношений $R_1$, $R_2$ содержит хотя бы один ненулевой шаг.
    Если для каждой точки $a\in\Int^2$ выполнено
    $$RC_{R_1}(a)=RC_{R_2}(a),$$
    то
    $$R_1 \approx R_2.$$
\end{namedLemma}

\textbf{Доказательство.}
Для $j\in\{1,2\}$ определим множество шагов отношения $R_j$:

$$S_{R_j}=\{v\in\Int^2\mid v\ne(0,0)\land R_j((0,0),v)\}.$$

По условию классы достижимости точки $(0,0)$ совпадают. Обозначим их через
$L$:

$$L=RC_{R_1}(0,0)=RC_{R_2}(0,0).$$

Тогда компоненту точки $(0,0)$ отношения $R_j$ можно рассматривать как граф
Кэли

$$\Cay(L,S_{R_j}).$$

По теореме~\ref{theorem:Cayley_Graphs_Isomorphisms} каждый автоморфизм этой
компоненты имеет вид $\varphi(x)=Ax+t$, где $t\in L$, $A(L)=L$ и
$A(S_{R_j})=S_{R_j}$.

Так как $S_{R_j}$ является конечным объединением множеств вида
$\{(\pm n,\pm k),(\pm k,\pm n)\}$, из условия
$A(S_{R_j})=S_{R_j}$ следует, что
$$A=
\begin{pmatrix}
  \pm1&0\\
  0&\pm1
\end{pmatrix}
\vee
A=
\begin{pmatrix}
  0&\pm1\\
  \pm1&0
\end{pmatrix}.$$

Компонентами обоих отношений являются одни и те же смежные классы по $L$.
На каждой компоненте автоморфизмы $R_1$ и $R_2$ имеют один и тот же вид,
описанный выше. Следовательно, группы автоморфизмов
$\langle\Int^2,R_1\rangle$ и $\langle\Int^2,R_2\rangle$ совпадают.

Рассмотрим общее полное вверх расширение этих структур. Его автоморфизмы
переставляют галактики и на каждой из них действуют автоморфизмом, описанным
выше. Поэтому каждый автоморфизм, сохраняющий $R_1$, сохраняет $R_2$, и
наоборот. Применяя теорему Свенониуса к обоим направлениям определимости,
получаем

$$R_1\approx R_2.$$

\hfill $\square$
\bs

\begin{theorem}
    \label{theorem:DnkUnionToDnk}
    $\forall N \in \Int\ \forall n_i, k_i \in \Int\ \exists n, k \in \Int\ \left(\bigvee \limits_{i = 1}^{N} D_{n_i, k_i}(x,y) \approx D_{n,k}(x,y)\right)$
\end{theorem}

\textbf{Доказательство.} Обозначим

$$G(x,y)=\bigvee\limits_{i=1}^{N}D_{n_i,k_i}(x,y).$$

Если все пары $(n_i,k_i)$ равны $(0,0)$, то $G=D_{0,0}$ и утверждение
доказано. Отношения $D_{0,0}$ можно удалить из остальных объединений, так как
они не меняют классы достижимости, а равенство определимо без параметров.

По доказанному выше критерию классов достижимости для каждого оставшегося
$D_{n_i,k_i}$ выберем отношение

$$C_i\in\{S_{d_i},D_{d_i}\}$$

с тем же классом достижимости точки $(0,0)$, где
$d_i=\gcd(n_i,k_i)$. Положим

$$G'(x,y)=\bigvee\limits_i C_i(x,y).$$

Тогда

$$RC_G(0,0)=RC_{G'}(0,0).$$

Все рассматриваемые отношения инвариантны относительно параллельных
переносов. Поэтому равенство классов точки $(0,0)$ влечёт равенство всех
классов достижимости. Остаётся найти отношение $S_d$ или $D_d$ с тем же
классом достижимости, что и у $G'$. Рассмотрим три случая.

\begin{enumerate}
    \item Все отношения $C_i$ имеют вид $S_{d_i}$.

    Тогда

    $$RC_G(0,0)=\sum\limits_i RC_{S_{d_i}}(0,0)
    =RC_{S_d}(0,0),$$

    где $d=\gcd(d_1,d_2,\dots,d_N)$. По
    лемме~\ref{lemm:RCtoR}

    $$G\approx S_d=D_{d,0}.$$

    \item Все отношения $C_i$ имеют вид $D_{d_i}$.

    Тогда

    $$RC_G(0,0)=\sum\limits_i RC_{D_{d_i}}(0,0)
    =RC_{D_d}(0,0),$$

    где $d=\gcd(d_1,d_2,\dots,d_N)$. По
    лемме~\ref{lemm:RCtoR}

    $$G\approx D_d=D_{d,d}.$$

    \item Среди отношений $C_i$ встречаются отношения обоих видов.

    Пусть

    $$a=\gcd\{d_i\mid C_i=S_{d_i}\},\qquad
    b=\gcd\{d_i\mid C_i=D_{d_i}\},$$

    и $d=\gcd(a,b)$. Запишем $a=dp$ и $b=dq$. Тогда
    $\gcd(p,q)=1$.

    Обозначим

    $$P=\{(u,v)\in\Int^2\mid u\equiv v\pmod 2\}.$$

    Поскольку

    $$RC_{S_a}(0,0)=a\Int^2,\qquad RC_{D_b}(0,0)=bP,$$

    получаем

    $$RC_G(0,0)=a\Int^2+bP=d(p\Int^2+qP).$$

    Рассмотрим два случая.

    \begin{enumerate}
        \item $p$ нечётно.

        Тогда $\gcd(p,2q)=1$. По лемме Безу существуют целые $u,v$, для
        которых

        $$up+2vq=1.$$

        Векторы $(p,0)$ и $(2q,0)=(q,q)+(q,-q)$ принадлежат
        $p\Int^2+qP$. Следовательно, $(1,0)$ и, аналогично, $(0,1)$ также
        принадлежат этому множеству. Поэтому

        $$p\Int^2+qP=\Int^2,$$

        и

        $$RC_G(0,0)=d\Int^2=RC_{S_d}(0,0).$$

        По лемме~\ref{lemm:RCtoR} $G\approx S_d.$

        \item $p$ чётно.

        Так как $\gcd(p,q)=1$, число $q$ нечётно. Поэтому

        $$p\Int^2+qP\subseteq P.$$

        По лемме Безу существуют целые $u,v$, для которых

        $$up+vq=1.$$

        Тогда

        $$(1,1)=u(p,p)+v(q,q),$$

        $$(1,-1)=u(p,-p)+v(q,-q).$$

        Значит, $P\subseteq p\Int^2+qP$, откуда

        $$p\Int^2+qP=P.$$

        Следовательно,

        $$RC_G(0,0)=dP=RC_{D_d}(0,0).$$

        По лемме~\ref{lemm:RCtoR} $G\approx D_d.$
    \end{enumerate}
\end{enumerate} \hfill $\square$

\begin{corollary}
    Из теорем \ref{theorem:DnkToSorD}, \ref{theorem:BinToDnk}, \ref{theorem:DnkUnionToDnk} следует, что любое определимое бинарное отношение эквивалентно $S_n$ или $D_n$.
 
\end{corollary}